\begin{ouverture}[thermo]{7}{Changements d’état et énergie}{Constitution et transformations}{Observer · Modéliser · Expérimenter · Expliquer}
\objectif{Modéliser un changement d’état à deux échelles.}
\objectif{Distinguer dissolution, fusion et transformation chimique.}
\objectif{Relier masse transformée et énergie échangée.}
\prerequis{Proportionnalité, conversions et lecture de graphiques.}
\end{ouverture}
\essentiel
\essSection{1}{Les six changements d’état}
Fusion : solide vers liquide ; solidification : liquide vers solide ; vaporisation : liquide vers gaz ; liquéfaction : gaz vers liquide ; sublimation : solide vers gaz ; condensation solide : gaz vers solide. L'espèce reste la même. L'écriture $\ce{H2O(s) -> H2O(l)}$ décrit une fusion.
\schema[corrige]{etats}
À l'échelle microscopique, les entités se réorganisent et leur mobilité change. Elles ne deviennent pas une autre espèce. La dissolution du sucre dans l'eau n'est pas une fusion : elle met en jeu un solvant.\schema[corrige]{etats-micro}
\essSection{2}{Température et état physique}
À pression fixée, un corps pur peut changer d'état à température constante : l'énergie reçue sert alors à modifier l'organisation de la matière. Un palier sur une courbe est un indice, à interpréter avec les conditions expérimentales. L'évaporation en surface peut avoir lieu sans ébullition de tout le liquide.
\schema[corrige]{palier}
\begin{attention}La masse totale se conserve dans un système fermé, mais le volume peut varier. L'eau solide est moins dense que l'eau liquide : une bouteille pleine peut être endommagée si elle gèle.\end{attention}
\essSection{3}{Le transfert thermique}
Une transformation endothermique reçoit de l'énergie thermique du milieu ; une transformation exothermique en cède. La fusion et la vaporisation sont endothermiques ; la solidification et la liquéfaction sont exothermiques.
\begin{loi}[Énergie de changement d'état]Pour une masse $m$ transformée à la température du changement d'état, la valeur de l'énergie échangée est $E=mL$, où $L$ est l'énergie massique correspondante en \unit{J.kg^{-1}}. Préciser séparément le sens du transfert.\end{loi}
\begin{exemple}Faire fondre \qty{0,100}{kg} de glace à 0 °C, avec $L_f=\qty{3,34e5}{J.kg^{-1}}$, demande \qty{3,34e4}{J}. Ce calcul ne couvre pas le chauffage préalable d'une glace à température négative.\end{exemple}
\essSection{4}{Expérimenter et modéliser}
\begin{experience}[Étudier un palier]Relever la température d'un corps pur pendant un chauffage doux, à intervalles réguliers ; noter les états présents. Tracer température en fonction du temps avec unités. Discuter les échanges avec l'environnement et la position de la sonde. Pour comparer l'énergie échangée à deux masses, garder le même changement d'état et les mêmes conditions.\end{experience}
\exercices
\rubrique{Prendre confiance}
\begin{blocLivre}{ex-7-1}\begin{exercice}[1]{Nommer les changements}
Des gouttes apparaissent sur une vitre froide ; un glaçon disparaît dans un verre.
\begin{questions}
\item Nommer les deux changements d’état.
\item L’espèce eau a-t-elle changé ?
\end{questions}
\end{exercice}\end{blocLivre}
\begin{blocLivre}{ex-7-2}\begin{exercice}[1]{Écrire symboliquement}
On transforme de l’eau liquide en vapeur.
\begin{questions}
\item Écrire l’équation du changement d’état.
\item Qualifier le transfert thermique pour l’eau.
\end{questions}
\end{exercice}\end{blocLivre}
\begin{blocLivre}{ex-7-3}\begin{exercice}[1]{Calculer une énergie de fusion}
On fait fondre 50,0 g de glace à 0 °C. $L_f=\qty{3,34e5}{J.kg^{-1}}$.
\begin{questions}
\item Convertir la masse en kilogrammes.
\item Calculer l’énergie reçue.
\end{questions}
\end{exercice}\end{blocLivre}
\begin{blocLivre}{ex-7-4}\begin{exercice}[1]{Dissoudre ou fondre ?}
Un sucre solide disparaît dans de l’eau à 20 °C ; un morceau de cire devient liquide quand on le chauffe.
\begin{questions}
\item Identifier les deux transformations.
\item Quelle opération implique un solvant ?
\end{questions}
\end{exercice}\end{blocLivre}
\rubrique{S’entraîner}
\begin{blocLivre}{ex-7-5}\begin{exercice}[2]{Lire une courbe thermique}
Une courbe idéalisée décrit le chauffage d’un corps pur à pression constante.
\schema{palier}
\begin{questions}
\item Identifier le changement pendant le palier.
\item Expliquer pourquoi la température n’augmente pas pendant cet intervalle.
\end{questions}
\end{exercice}\end{blocLivre}
\begin{blocLivre}{ex-7-6}\begin{exercice}[2]{Une masse à solidifier}
La solidification libère 66,8 kJ. $L_f=\qty{334}{kJ.kg^{-1}}$.
\begin{questions}
\item Déterminer la masse d’eau solidifiée.
\item Quel milieu reçoit l’énergie ?
\end{questions}
\end{exercice}\end{blocLivre}
\begin{blocLivre}{ex-7-7}\begin{exercice}[2]{Deux masses de glace}
Deux blocs de même glace à 0 °C ont des masses de 40 g et 100 g.
\begin{questions}
\item Comparer les énergies nécessaires à leur fusion complète.
\item Faut-il utiliser la même valeur de $L_f$ ?
\end{questions}
\end{exercice}\end{blocLivre}
\begin{blocLivre}{ex-7-8}\begin{exercice}[2]{Une évaporation qui refroidit}
De l’eau s’évapore sur la peau.
\begin{questions}
\item Dans quel sens l’énergie est-elle transférée ?
\item Pourquoi la peau peut-elle refroidir ?
\end{questions}
\end{exercice}\end{blocLivre}
\rubrique{Raisonner et approfondir}
\begin{blocLivre}{ex-7-9}\begin{exercice}[3]{Une puissance utile fournie}
Un dispositif fournit effectivement \qty{100}{J} chaque seconde à de la glace à 0 °C. La masse est 0,200 kg et $L_f=\qty{3,34e5}{J.kg^{-1}}$.
\begin{questions}
\item Calculer l’énergie de fusion.
\item Estimer la durée de fusion, sans autre perte.
\end{questions}
\end{exercice}\end{blocLivre}
\begin{blocLivre}{ex-7-10}\begin{exercice}[3]{Masse conservée, volume changé}
100 g d’eau liquide deviennent de la glace. $\rho_l=\qty{1,00}{g.cm^{-3}}$, $\rho_s=\qty{0,92}{g.cm^{-3}}$.
\begin{questions}
\item Calculer les volumes avant et après.
\item Pourquoi un récipient rigide plein pose-t-il problème ?
\end{questions}
\end{exercice}\end{blocLivre}
\begin{blocLivre}{ex-7-11}\begin{exercice}[3]{Une étape oubliée}
On veut transformer de la glace à -10 °C en eau à 20 °C.
\begin{questions}
\item Pourquoi $E=mL_f$ seul ne suffit-il pas ?
\item Lister les trois étapes physiques.
\end{questions}
\end{exercice}\end{blocLivre}
\begin{blocLivre}{ex-7-12}\begin{exercice}[3]{Comparer deux matériaux}
Pour la même masse, A a une énergie de fusion massique de 180 kJ/kg et B de 300 kJ/kg.
\begin{questions}
\item Lequel stocke davantage d’énergie par fusion ?
\item Quelle autre condition faut-il vérifier pour un usage à 20 °C ?
\end{questions}
\end{exercice}\end{blocLivre}

